\section{Attach the binary scale to the value}
A stored Q15 coefficient \(c\) represents \(c/2^{15}\). For signed halfword codes,
its range is \(-1\) through \(32767/32768\). Positive one has no single signed
halfword encoding at this scale. Two coefficients 16384 each represent a total
weight of one; eight coefficients 4096 each do too.

Input PCM codes remain signed integers for the kernel. Multiplying code \(u\) by
coefficient code \(c\) gives \(uc\), whose conversion back to sample-code units
requires division by 32768. If input codes are also viewed as normalized values
\(u/32768\), products have thirty fractional bits. These are two descriptions of
the same integer calculation, not different instructions or stored types.

The supplied moving average uses eight coefficients 4096. Each output is therefore
\(\sum u/8\), with the specified final rounding. Although every individual weight
is small, the general interface permits every signed halfword coefficient. Correctness
cannot depend on the default coefficients or a presumed gain no greater than one.

\begin{table}[htbp]
  \centering
  \caption{Stored coefficient codes and their exact weights.}
  \tagpdfsetup{table/header-rows={1}}
  \begin{tabular}{rr}
    \toprule
    Code   & Weight        \\
    \midrule
    -32768 & \(-1\)        \\
    -16384 & \(-1/2\)      \\
    0      & \(0\)         \\
    4096   & \(1/8\)       \\
    16384  & \(1/2\)       \\
    32767  & \(1-1/32768\) \\
    \bottomrule
  \end{tabular}
\end{table}
